% T31 source: Примеры базовой модели.tex lines 628--821
\detailedexampleheading{Variant 3. Delegated Agents with Detailed Specifications}

In this variant, four AI agents work in parallel ($P = 4$). The configuration
is specified as
$\sigma_3 = (P{=}4,\ \text{delegated mode with detailed specifications})$.
The coefficients $k_i$ are lower than in the preceding variants because the
operating-mode profile changes, not because parallelism increases: the agents
receive detailed specifications and execute the autonomous agent phases, after
which the developer accepts the result. The dependency graph among the
aggregate tasks remains unchanged and limits the exploitable parallelism.

\paragraph{Parameters.}
\begin{table}[ht!]
\centering
\caption{Parameters of Variant 3}
\label{tab:practical-v3-params}
\begin{tabular}{clccc}
\toprule
$i$ & Task & $Z_i$ & $k_i$ & $h_i$ \\
\midrule
1 & API contract and request validation & 2 & 0.65 & 0.25 \\
2 & Persistence layer & 3 & 0.70 & 0.25 \\
3 & Business logic & 4 & 0.80 & 0.25 \\
4 & Background worker & 3 & 0.75 & 0.25 \\
5 & Observability and configuration & 2 & 0.60 & 0.25 \\
6 & Tests & 3 & 0.65 & 0.25 \\
7 & Packaging and deployment & 2 & 0.70 & 0.25 \\
\bottomrule
\end{tabular}
\end{table}

The same human component, $h_i = 0.25$, is assumed for all tasks. The
developer's involvement consists of preparing a detailed specification and
subsequently reviewing the result; implementation is delegated entirely to the
agent.

\paragraph{Calculations.}
\begin{align*}
  \sum_{i=1}^{7} Z_i k_i &= 2{\cdot}0.65 + 3{\cdot}0.70 + 4{\cdot}0.80 + 3{\cdot}0.75 + 2{\cdot}0.60 + 3{\cdot}0.65 + 2{\cdot}0.70 \\
  &= 1.30 + 2.10 + 3.20 + 2.25 + 1.20 + 1.95 + 1.40 = 13.40.
\end{align*}

The weighted-average normalized duration is
\[
  \bar{k}_w = \frac{13.40}{19} \approx 0.705.
\]

The task weights in the graph are $w_i = Z_i k_i X$:
\[
  \{w_i\}_{i=1}^{7} = \{5.2,\; 8.4,\; 12.8,\; 9.0,\; 4.8,\; 7.8,\; 5.6\}~\text{h}.
\]

The total work is
\[
  W = \sum_{i=1}^{7} w_i = 13.40 \cdot 4 = 53.6~\text{h}.
\]

The idealized completion time for independent tasks, that is, the lower bound
based on total work when $P = 4$, is
\[
  T_{ai}^{(0)} = \frac{W}{P} = \frac{53.6}{4} = 13.4~\text{h}.
\]

The idealized speedup without accounting for dependencies is
\[
  S_E = \frac{76}{13.4} \approx 5.67.
\]

The critical path through the graph comprises Tasks
$1 \to 2 \to 3 \to 4 \to 6$:
\[
  L = 5.2 + 8.4 + 12.8 + 9.0 + 7.8 = 43.2~\text{h}.
\]

The human time at level M4 is
\[
  \sum_{i=1}^{7} Z_i h_i = 0.25 \cdot 19 = 4.75,
  \qquad
  H = 4.75 \cdot 4 = 19.0~\text{h},
  \qquad
  \bar{h}_w = 0.25,
  \qquad
  \frac{1}{\bar{h}_w} = 4.0.
\]

The lower bound on the makespan accounting for the dependency graph and the
human resource is
\[
  B_4 = \max\left\{\frac{W}{P},\; L,\; H\right\} = \max\{13.4,\; 43.2,\; 19.0\} = 43.2~\text{h}.
\]

One feasible assignment to four agents that respects the dependencies is
\begin{itemize}
  \item Agent 1: Task 1 $[0,5.2]$, Task 2 $[5.2,13.6]$, Task 3
  $[13.6,26.4]$, Task 4 $[26.4,35.4]$, and Task 6 $[35.4,43.2]$.
  \item Agent 2: Task 5 $[6.7,11.5]$ and Task 7 $[15.6,21.2]$.
  \item Agents 3 and 4 are idle because, at this coarse graph granularity, the
  critical path dominates the exploitable parallelism.
\end{itemize}

This schedule attains the lower bound:
\[
  T(\pi_3)=T^{*}=B_4=43.2~\text{h}.
\]

\paragraph{Human-Resource Check (Level M4).}
To prevent an agent from starting work before task specification, the human
time is divided equally between preparation of the input specification and
final review; the agent performs the implementation between these stages. The
developer's blocks are scheduled as follows: Task 1---$[0,\,1.0]$ and
$[4.2,\,5.2]$; Task 2---$[5.2,\,6.7]$ and $[12.1,\,13.6]$; Task
5---$[6.7,\,7.7]$ and $[10.5,\,11.5]$; Task 3---$[13.6,\,15.6]$ and
$[24.4,\,26.4]$; Task 7---$[15.6,\,16.6]$ and $[20.2,\,21.2]$; Task
4---$[26.4,\,27.9]$ and $[33.9,\,35.4]$; and Task 6---$[35.4,\,36.9]$
and $[41.7,\,43.2]$. The intervals are pairwise nonoverlapping, and each input
specification precedes the agent component of the corresponding task. Total
human utilization is $19.0$~h $= H$ over a project duration of $43.2$~h, so the
schedule is feasible without modification.

The achievable speedup accounting for the graph and the human resource is
\[
  S_E^{\text{ach}} = \frac{76}{43.2} \approx 1.76.
\]

\paragraph{Interpretation.}
With four agents, total work decreases to $W=53.6$~h, and the estimate for
independent tasks is $W/P=13.4$~h. Under the coarsened graph, however, this
value is unattainable: the chain $1 \to 2 \to 3 \to 4 \to 6$ has length
$43.2$~h and becomes the primary constraint. After dependencies are taken into
account, the speedup is therefore not $\times 5.67$ but $\times 1.76$.

The improvement relative to Variant~2 should be interpreted with caution. The
critical path decreases from $47.8$ to $43.2$~h because the coefficients $k_i$
decrease upon switching to delegated mode with detailed specifications. The
increase in $P$ from 2 to 4 does not itself affect $L$: additional agents do not
accelerate the sequential chain of aggregate tasks and are idle in this
schedule. The next change to the scenario is therefore not a further increase
in the number of agents, but the decomposition of large tasks on the critical
path into smaller, independently verifiable components. Variant~4 considers
precisely this change.

Level M4 imposes an additional upper bound on speedup. In this variant, the
human-resource ceiling $1/\bar{h}_w = 4.0$ is substantially higher than the
achieved speedup of $\times 1.76$, so $L$ remains the active constraint.
Nevertheless, under the selected mode ($\bar{h}_w = 0.25$), neither
decomposition nor any number of agents can provide speedup greater than
$\times 4.0$: the 19 hours of sequential human time cannot be eliminated. As
the critical path is shortened, $H$ becomes the active constraint.

\detailedexampleheading{Variant 4. The Same Configuration with Finer Decomposition}

This variant differs from Variant~3 only in its decomposition. The number of
agents $P = 4$, delegated mode with detailed specifications, the coefficients
$k$, and the human components $h$ remain unchanged. Only the graph changes: the
aggregate Task 6, ``Tests,'' is replaced by the three subtasks from
Table~\ref{tab:practical-subtasks}, each with its exact dependencies. In
Variants 1--3, the configuration changed while the decomposition remained
fixed; here, by contrast, the decomposition changes while the configuration
remains fixed. The difference from Variant~3 is therefore attributable solely
to decomposition.

\paragraph{Parameters.}
Tasks 1--5 and 7 remain unchanged (Table~\ref{tab:practical-v3-params}). Task 6
is decomposed as specified in Table~\ref{tab:practical-subtasks}:

\begin{table}[ht!]
\centering
\footnotesize
\caption{Subtasks of Task 6 in Variant 4}
\label{tab:practical-v4-params}
\begin{tabular}{clcccl}
\toprule
$i$ & Subtask & SP & $Z_i$ & $k_i$ / $h_i$ & Depends on \\
\midrule
6a & Unit tests for the business logic & 2 & 1 & 0.65 / 0.25 & 3 \\
6b & Integration tests for the API and persistence layer & 3 & 1.5 & 0.65 / 0.25 & 1, 2 \\
6c & Verification of background processing and retries & 1 & 0.5 & 0.65 / 0.25 & 4 \\
\bottomrule
\end{tabular}
\end{table}

The decomposition is assumed \emph{not to change the amount of work}: the
subtasks inherit $k = 0.65$ and $h = 0.25$ from the original task. The
quantities $W$, $\bar{k}_w$, and $H$ are therefore unchanged, so the comparison
with Variant~3 isolates the structural effect. In practice, decomposition
entails additional overhead for preparing specifications, issuing tasks, and
accepting each subtask. Under the adopted accounting rule, this overhead is
included in $k^{(r)}$ and increases it when decomposition becomes excessive.
Conversely, $k$ may decrease for small tasks with clear boundaries. Thus, holding $k$ and
$h$ constant is a neutral scenario assumption, and the sensitivity of the
result to decomposition overhead requires separate calibration.

The subtask sizes are fractional ($Z_{6\text{b}} = 1.5$,
$Z_{6\text{c}} = 0.5$): an odd number of story points under the conversion
$2$~SP $= 1\,Z$. The subtask weights are
\[
  w_{6\text{a}} = 1 \cdot 0.65 \cdot 4 = 2.6~\text{h}, \qquad
  w_{6\text{b}} = 1.5 \cdot 0.65 \cdot 4 = 3.9~\text{h}, \qquad
  w_{6\text{c}} = 0.5 \cdot 0.65 \cdot 4 = 1.3~\text{h};
\]
together, they total $7.8$~h, the weight of the original Task 6.

\paragraph{Calculations.}
The aggregate quantities are the same as in Variant~3:
\[
  W = 53.6~\text{h}, \qquad
  \bar{k}_w \approx 0.705, \qquad
  \frac{W}{P} = 13.4~\text{h}, \qquad
  H = 19.0~\text{h}, \qquad
  \frac{1}{\bar{h}_w} = 4.0.
\]

The critical path changes. The monolithic vertex 6 and its incoming edges
$3 \to 6$ and $4 \to 6$ are replaced by three vertices with the exact
dependencies $3 \to 6\text{a}$, $\{1,2\} \to 6\text{b}$, and
$4 \to 6\text{c}$, and the longest path becomes
\[
  1 \to 2 \to 3 \to 4 \to 6\text{c}: \qquad
  L = 5.2 + 8.4 + 12.8 + 9.0 + 1.3 = 36.7~\text{h}
\]
(the competing paths are shorter: $1 \to 2 \to 3 \to 6\text{a}$ yields
$29.0$~h, $1 \to 2 \to 6\text{b}$ yields $17.5$~h, and
$1 \to 5 \to 7$ yields $15.6$~h). The tail of the critical path decreases
from $7.8$~h (all of Task 6) to $1.3$~h (only the component that actually waits
for the background worker).

The lower bound on the makespan is
\[
  B_4 = \max\left\{\frac{W}{P},\; L,\; H\right\} = \max\{13.4,\; 36.7,\; 19.0\} = 36.7~\text{h}.
\]

A feasible assignment that respects the dependencies is
\begin{itemize}
  \item Agent 1: Task 1 $[0,\,5.2]$, Task 2 $[5.2,\,13.6]$, Task 3
  $[13.6,\,26.4]$, Task 4 $[26.4,\,35.4]$, and Task 6c
  $[35.4,\,36.7]$.
  \item Agent 2: Task 5 $[6.7,\,11.5]$; for Task 7, the input specification
  occupies $[19.5,\,20.5]$, agent implementation occupies
  $[20.5,\,24.1]$, and the result then waits in the queue for developer
  attention until the review at $[28.4,\,29.4]$. During the waiting interval
  $[24.1,\,28.4]$, the agent remains assigned to Task 7 and cannot receive
  other work.
  \item Agent 3: Task 6b $[15.6,\,19.5]$ and Task 6a
  $[27.9,\,30.5]$. Agent 4 is idle.
\end{itemize}

This schedule attains the lower bound:
\[
  T(\pi_4)=T^{*}=B_4=36.7~\text{h}.
\]

\paragraph{Human-Resource Check (Level M4).}
As in Variant~3, the human time for each task is divided equally between the
input specification and final review. The developer's blocks are as follows:
Task 1---$[0,\,1.0]$ and $[4.2,\,5.2]$; Task 2---$[5.2,\,6.7]$ and
$[12.1,\,13.6]$; Task 5---$[6.7,\,7.7]$ and $[10.5,\,11.5]$; Task
3---$[13.6,\,15.6]$ and $[24.4,\,26.4]$; Task 6b---$[15.6,\,16.35]$
and $[18.75,\,19.5]$; Task 7---$[19.5,\,20.5]$ and $[28.4,\,29.4]$;
Task 4---$[26.4,\,27.9]$ and $[33.9,\,35.4]$; Task
6a---$[27.9,\,28.4]$ and $[30.0,\,30.5]$; and Task
6c---$[35.4,\,35.65]$ and $[36.45,\,36.7]$. The intervals are pairwise
nonoverlapping, each input specification precedes agent implementation, and
total human utilization is $19.0$~h $= H$. The schedule is therefore feasible
with respect to the human resource and does indeed attain $36.7$~h.

The achievable speedup accounting for the graph and the human resource is
\[
  S_E^{\text{ach}} = \frac{76}{36.7} \approx 2.07.
\]

\paragraph{Interpretation.}
The makespan decreases from $43.2$ to $36.7$~h, while speedup increases from
$\times 1.76$ to $\times 2.07$, an increase of $18\%$, without adding agents,
changing the operating mode, or improving the tool. This effect is obtained by
replacing one aggregate vertex with three subtasks already identified in
Table~\ref{tab:practical-subtasks}. The dependency ``testing after completion
of the entire core'' was an artifact of graph coarsening; refining it shortens
the critical path by $6.5$~h. In the comparison considered here, the result is
determined by decomposition, not by the number of agents.

Only the transition from Variant 3 to Variant 4 is an all-else-equal
comparison: the decrease from $43.2$ to $36.7$~h is attributable to
decomposition with $P$, $k$, and $h$ held constant. In the transition from
Variant 2 to Variant 3, both $P$ and the operating mode change, so their effects
are not separated. The calculation nevertheless shows that the reduction in
the active bound $L$ from $47.8$ to $43.2$~h is associated with lower $k$
values along the critical path. The value of $P$ does not itself enter the
formula for $L$, and the additional agents are idle.

The comparative conclusion for Variants 3 and 4 also holds without an exact
absolute calibration. By Proposition~\ref{prop:invariance}, multiplying all
agent and human-attention phases in both variants by a common factor
$\kappa>0$ multiplies both makespans by $\kappa$ but leaves the ratio
$43.2/36.7$ and the ranking of the variants unchanged. This robustness does
not extend to task-specific errors of unequal magnitude or to changes in the
relative profiles of $k_i$ and $h_i$.

Decomposition can be continued iteratively. The next candidate is the
dependency $3 \to 4$, which at the subtask level reduces to the 2-SP state model
within the 8-SP task. If it is separated into an early subtask, the background
worker can begin before the remaining business-logic use cases are complete.
The limit to such decomposition is already apparent: the path
$1 \to 2 \to 3 \to 4$, of length $35.4$~h, almost completely determines
$L = 36.7$~h, while the makespan is bounded below by the human resource
$H = 19.0$~h and speedup is bounded above by $\times 4.0$. As the critical path
is shortened further, $H$ becomes the active constraint, as predicted by level
M4.
